We observe a binary system, at a distance $d = 89$\,pc, with a circular,
edge-on orbit around the common center of mass and period of 100 days. Our
space-based telescope is of diameter $D = 10$\,m and operates at a wavelength
$\lambda = 364$\,nm. We notice that for a total of 38 days during each full
period, the two stars cannot be resolved by this telescope as separate
objects. The wavelength of peak emission for star A is
$\lambda_A = 500$\,nm and that for star B is $\lambda_B = 600$\,nm.

\begin{description}
\item[(a)] What are the temperatures of the stars A and B, ($T_A$, $T_B$)?
  \hfill(2pt)
\item[(b)] Calculate the distance $l$ between the stars. \hfill(6pt)
\item[(c)] Calculate the sum of the masses of the stars ($M_T$). \hfill(2pt)
\end{description}

Combined photometry of the system:

\begin{center}
\begin{tabular}{clcccc}
 & Configuration & $U_0$ & $(U-B)$ & $(B-V)$ & BC \\
1 & Stars next to each other & 6.39 & 0.2 & 0.1 & 0.1 \\
2 & B transiting in front of A & 6.86 & 0.25 & 0.12 & 0.175 \\
\end{tabular}
\end{center}

The interstellar extinction per kpc of $U$ is $a_U = 1.4$\,mag/kpc. The
ratio of the densities of the stars is $\rho_A/\rho_B = 0.7$.

Note: for bolometric correction we use the convention:
\[ \mathrm{BC} = m_{\mathrm{bol}} - m_V \]

\begin{description}
\item[(d)] Calculate the masses of both the stars ($M_A$, $M_B$).
  \hfill(10pt)
\end{description}
